\subsection{Volumetric Encoders and Attention}\label{sec:fundamentals-backbone}

\paragraph{Convolutional volumetric backbones.}
The 3D U-Net~\cite{cciccek20163d} architecture encodes a volume into a feature pyramid with a
downsampling encoder and decodes it back to full resolution. Deeper stages have lower spatial resolution and more channels, while shallow stages keep spatial detail at a low channel count. Our model uses this split directly: the coarse, high-channel stages feed the attention layers, and the fine stages are passed to the decoder as skip connections (\autoref{sec:methodology-overview}). UNETR~\cite{hatamizadeh2022unetr} is an earlier example of combining
transformer attention with a convolutional decoder for 3D segmentation.

\paragraph{Attention.}
Scaled dot-product attention~\cite{vaswani2017attention} computes, for each
query, a weighted average of value vectors. The weights are given by the
similarity between the query and each key:
%
\begin{equation}
  \mathrm{Attn}(Q,K,V) = \mathrm{softmax}\!\left(\frac{QK^\top}{\sqrt{d}}\right)V,
  \label{eq:attention}
\end{equation}
%
where $d$ is the key dimension. Multi-head attention runs several attention
operations in parallel on lower-dimensional projections and concatenates their
outputs, so that different heads can capture different relations within the
same input.

\paragraph{Position encoding.}
Attention is permutation-invariant, so positional information has to be added
explicitly. Rotary position embedding (RoPE)~\cite{su2021roformer} rotates each
query and key vector by an angle proportional to its position before the dot
product. The resulting attention score depends only on the relative position
of the two tokens. In 3D, RoPE can be applied per axis: the head dimension is
split into three parts, and each part is rotated according to the token's
coordinate along one spatial axis. If these coordinates are given in physical
units, i.e.\ scaled by the voxel spacing, attention between tokens from volumes
of different resolution reflects their physical distance rather than their grid
indices. Our coarse-to-fine approach relies on this property, since its levels operate at different spacings (\autoref{sec:methodology-cascade}).